\section{Diseño del control supervisor (Nivel 1)}

El control supervisor actúa como el cerebro lógico-operativo de la grúa. Se ejecuta a un
tiempo de muestreo $T_{s1} = 20$ ms y su objetivo es tomar decisiones lógicas, coordinar la
operación de ambos ejes y garantizar la seguridad funcional del sistema, habilitando o
bloqueando las consignas que luego procesa el control regulatorio. A diferencia del
Nivel 2, que opera con variables continuas muestreadas, el Nivel 1 es un sistema de eventos
discretos (DEDS) implementado como una máquina de estados en Stateflow.

Este nivel recibe como entradas la realimentación de los sensores de la planta (posiciones de
encoder $x_{t}$ e $y_{h}$, velocidades estimadas por los observadores, tensión de cable $F_{hw}$, ángulo de
balanceo $\theta_l$), las señales del joystick del operador, los finales de carrera y las
señales de falla provenientes del Nivel 0. Como salidas genera las consignas de velocidad
$v_{x\ suave}$ y $v_{y\ suave}$, los
comandos de freno de operación \textcolor{signalcolor}{$BRK_t$} y \textcolor{signalcolor}{$BRK_h$},
la señal de twistlocks \textcolor{signalcolor}{TLK}, el \textit{heartbeat} y diversas
señales de estado de la interfaz. 
Para organizar esta lógica, la arquitectura se divide en múltiples procesos paralelos
que el autómata evalúa simultáneamente en cada paso.

\subsection{Autómata principal de modos de operación}

Como bloque principal se desarrolla la máquina de estados \texttt{AUTOMATA\_PRINCIPAL}. Este es
el hilo de ejecución primario que dicta el comportamiento global de la grúa, gestionando de
manera segura las transiciones entre los distintos modos de operación.

\begin{itemize}
    \item \textbf{Reposo (\texttt{REPOSO}):} estado inicial y de espera de comandos. La grúa
    permanece detenida con los frenos cerrados. Es el estado al que se regresa cuando finaliza el movimiento.

    \item \textbf{Manual (\texttt{MANUAL}):} permite al operador tomar el control directo de las
    velocidades de ambos ejes mediante el joystick. Las consignas del joystick se traducen a
    referencias de velocidad para el control regulatorio, permitiendo movimientos precisos y
    suaves. Se activa al detectarse una consigna en el joystick.

    \item \textbf{Modo automático (\texttt{MODO\_AUTO}):} es el modo de operación principal del
    sistema. Una vez que el operador selecciona las posiciones de origen y destino, el autómata
    gestiona de forma autónoma el ciclo completo de traslado del contenedor: planificación de
    alturas seguras, movimiento coordinado de ambos ejes y operación de los twistlocks. La misión
    se organiza en tres etapas consecutivas, identificadas internamente por la variable
    \texttt{estado}: la etapa de recolección (\texttt{estado}~$= 1$), la etapa de depósito
    (\texttt{estado}~$= 2$) y el retorno a la posición de inicio. Los subestados que implementan
    este ciclo son:

    \begin{itemize}
        \item \texttt{AUTO\_INIT}: inicializa la maniobra. Guarda la posición horizontal actual
        del carro como punto de retorno al final de la misión (\textit{home}) y determina la
        etapa de inicio según el estado de los twistlocks: si el spreader ya tiene carga
        enganchada (\textcolor{signalcolor}{TLK}~$= 1$), se salta directamente a la etapa de
        depósito.

        \item \texttt{SEL\_ORIGEN} y \texttt{SEL\_DESTINO}: permiten al operador navegar el
        perfil de obstáculos discretizado con el joystick para seleccionar los índices de
        columna de origen y destino del contenedor. La selección incorpora un filtro de
        antirrebote de aproximadamente 500~ms para evitar cambios involuntarios. El sistema
        valida que exista un contenedor en la columna de origen y que el slot de destino esté
        disponible; al confirmar con el botón \textcolor{signalcolor}{BTN\_TLK} se fija el
        índice correspondiente y se avanza al subestado siguiente.

        \item \texttt{SELECCION\_TRAYECTORIA}: según la etapa activa del ciclo (ir a buscar, ir
        a depositar o retornar), selecciona el índice objetivo y calcula la altura mínima de
        tránsito seguro entre la posición actual del carro y dicho objetivo:
        \begin{equation}
            y_{\text{segura}} = \max_{k \,\in\, [i_{\min},\, i_{\max}]} y_{c0,\text{vec}}[k]
            + H_c + 1\,\text{m}
            \label{eq:altura_segura}
        \end{equation}
        donde $[i_{\min}, i_{\max}]$ es el rango de columnas que el carro atravesará. El margen
        $H_c + 1$~m garantiza que el spreader, con o sin contenedor, despeje cualquier
        obstáculo intermedio.

        \item \texttt{MOVIMIENTO\_CURVO} y \texttt{MOVIMIENTO\_CURVO\_FINAL}: ejecutan el
        desplazamiento coordinado en dos fases. Primero, el izaje lleva el spreader hasta
        $y_{\text{segura}}$ (Ec.~\ref{eq:altura_segura}) antes de iniciar el
        desplazamiento horizontal; a continuación, el carro traslada el spreader hasta el
        índice objetivo. Ambos ejes utilizan un controlador proporcional-derivativo de
        seguimiento de posición que genera las consignas de velocidad $v_{x\,target}$ y
        $v_{y\,target}$:
        \begin{equation}
            v_{target}(k) = K_{p} \, \bigl(pos_{target}(k) - pos_{med}(k)\bigr) - K_{d} \, v_{obs}(k)
            \label{eq:control_pd}
        \end{equation}
        siendo $pos_{med}(k)$ la posición medida por el encoder y $v_{obs}(k)$ la velocidad
        estimada por el observador de Luenberger del Nivel~2. Durante toda esta fase, el
        autómata monitorea la consigna del joystick: si el operador realiza una entrada
        voluntaria, el sistema transfiere el control a \texttt{MODO\_MANUAL\_SIMPLE} para
        que el operador complete el posicionamiento fino, con la posibilidad de retomar la
        secuencia automática.

        \item \texttt{MODO\_MANUAL\_SIMPLE}: etapa de aproximación final con o sin contenedor
        realizada por el operador a velocidad reducida (hasta un 5\,\% de $v_{t\,\max}$ en
        horizontal), que permite el alineamiento preciso del spreader sobre el contenedor o
        el slot de destino. Al confirmar con \textcolor{signalcolor}{BTN\_TLK}, el autómata
        cierra o abre los twistlocks según corresponda y avanza a la etapa siguiente del ciclo.
    \end{itemize}

    La velocidad máxima de izaje durante la etapa de depósito se ajusta dinámicamente según
    la curva de potencia constante del motor (Ec.~\ref{eq:curva_potencia_constante}), en función
    de la masa del contenedor estimada por el filtro de la celda de carga en el subestado
    \texttt{CARGADO} del proceso de twistlocks.

    \item \textbf{Estados de error:} ante una falla propia del supervisor o detectada por el
    Nivel 0, el autómata aborta la operación en curso y fuerza al sistema a un estado seguro con
    los frenos aplicados. Los errores se discriminan según su origen: \texttt{ERROR\_FC\_sup\_inf}
    para fallas de finales de carrera de izaje, \texttt{ERROR\_FC\_izq\_der} para los del carro y
    \texttt{ERROR\_OTRO} para el resto de las condiciones críticas, cerrando los frenos correspondientes según el error.


    \item \textbf{Recuperación (\texttt{RECUPERACION}):} estado de resguardo que limita la
    velocidad y dirección del movimiento manual para que el operador pueda llevar la grúa a una
    condición segura tras una falla. Se activa con el botón de reset después de una falla y se
    abandona cuando la posición se aleja del fin de carrera correspondiente, si no hay uno se abandona inmediatamente.


\end{itemize}

\begin{figure}[H]
    \centering
    \includegraphics[width=0.75\textwidth]{nivel_1/automata_principal.png}
    \caption{Máquina de estados del autómata principal en Stateflow}
    \label{fig:n1_automata_principal}
\end{figure}

\subsection{Gestión de twistlocks}

Este proceso paralelo se encarga exclusivamente de la lógica del spreader,
administrando las confirmaciones de los sensores para tomar o soltar el
contenedor, asegurando que la grúa nunca intente un movimiento de izaje si el anclaje no es
mecánicamente seguro. Posee cuatro estados internos:

\begin{itemize}
    \item \texttt{VACIO}: estado inicial, con el spreader sin carga moviéndose a la velocidad
    máxima permitida.
    \item \texttt{CERRANDO\_TLK}: se activa cuando la tensión en el cable disminuye lo suficiente
    para indicar que el spreader se apoyó sobre un contenedor, se confirma que hay un contenedor y no es el piso y 
    se presiona el botón TLK. En este estado
    se cierran los twistlocks
    (\textcolor{signalcolor}{TLK} $= 1$) y se limita la velocidad de izaje.
    \item \texttt{CARGADO}: con el contenedor enganchado, se estima su masa a partir de la celda
    de carga mediante un filtro de primer orden sobre la tensión del cable. Con dicha estimación
    se ajusta dinámicamente el límite de velocidad de izaje según la curva de potencia constante
    del motor (Ec. \ref{eq:curva_potencia_constante}) y se detecta una condición de sobrecarga si la masa estimada supera el máximo
    admisible.
    \begin{equation}
        v_{h,lim}({m}_l) = 
        \begin{cases} 
            v_{h\max} & \text{si } {m}_l < 32500 kg \\
            \frac{P_{h\ nom}}{{m}_l \, g} & \text{si } {m}_l \ge 32500 kg
        \end{cases}
        \label{eq:curva_potencia_constante}
    \end{equation}

    \item \texttt{ABRIENDO\_TLK}: de forma análoga al cierre, detecta la disminución de tensión
    al apoyar la carga para liberar los twistlocks y espera la pulsación del botón TLK. Una vez superada cierta altura de seguridad,
    se transiciona nuevamente al estado inicial.
\end{itemize}


\begin{figure}[H]
    \centering
    \includegraphics[width=0.65\textwidth]{nivel_1/twistlocks.png}
    \caption{Lógica de gestión de twistlocks en Stateflow}
    \label{fig:n1_twistlocks}
\end{figure}

\subsection{Suavizador de consignas de velocidad}

Este proceso aplica filtros de suavizado a las consignas de velocidad provenientes del autómata
principal. Resulta especialmente importante en el modo manual, donde las señales del joystick
pueden ser abruptas. El bloque implementa límites de \textit{jerk} y aceleración, garantizando
que cualquier cambio en la referencia se realice de manera gradual y evitando movimientos
bruscos que comprometan la estabilidad o la seguridad de la operación. Las consignas crudas
$v_{x\ target}, v_{y\ target}$ se transforman en
consignas suavizadas $v_{x\ suave}, v_{y\ suave}$.

\begin{figure}[H]
    \centering
    \includegraphics[width=0.6\textwidth]{nivel_1/suavizado_velocidad.png}
    \caption{Suavizador de consignas de velocidad en Stateflow}
    \label{fig:n1_suavizado}
\end{figure}

\subsection{Análisis espacial, perfil de obstáculos e índices}

Con el fin de que el modo automático pueda trazar rutas seguras, la grúa debe conocer su entorno y tener un mapa de alturas o perfil de obstáculos
del mismo.  Para simplificar la programación de este proceso,
el aútomata conoce de antemano el perfil inicial y se define como condición de habilitación del modo automático 
que el perfil de obstáculos haya sido previamente escaneado en su totalidad.
Además, el sistema se encarga de cambiar el perfil acorde los contenedores que se van moviendo. 

Al mismo tiempo, calcula el índice discreto en el que se encuentra el carro, y genera una envolvente de seguridad a partir
del perfil de obstáculos, con márgenes laterales y superiores que garantizan que la grúa no intente
pasar por zonas donde pudiera colisionar. Esto es usado en la planificación de trayectorias para generar consignas de velocidad seguras.

\begin{figure}[H]
    \centering
    \includegraphics[width=0.46\textwidth]{nivel_1/indices_perfil.png}
    \caption{Cálculo de índice y actualización del perfil de obstáculos en Stateflow}
    \label{fig:n1_indices_perfil}
\end{figure}
\subsection{Límites suaves de velocidad}

Este proceso actúa como una barrera de seguridad preventiva. Antes de que el carro o el izaje
alcancen sus respectivos fines de carrera, este bloque recorta gradualmente las referencias de velocidad, forzando una detención
suave y controlada dentro del área de trabajo válida.

El recorte se aplica de forma independiente en los cuatro límites de operación. Para cada uno se calcula una velocidad segura
de aproximación a partir de la distancia remanente al límite, de modo que la velocidad
disponible se reduce a medida que la carga se acerca al borde, considerando además un retardo
dinámico que anticipa el frenado en función de la velocidad actual. Si la consigna suavizada
excede la velocidad segura calculada, se satura a esta última, garantizando que el sistema nunca
llegue a los finales de carrera a velocidad elevada.
\newpage
%\thispagestyle{empty}
\begin{figure}[H]
    \centering
    \includegraphics[width=0.5\textwidth]{nivel_1/limites_suaves.png}
    \caption{Lógica de límites suaves de velocidad en Stateflow}
    \label{fig:n1_limites_suaves}
\end{figure}
\newpage

\subsection{Supervisión de comunicación (Watchdog Heartbeat)}

Este proceso constituye un anillo de seguridad vital. Se encarga de emitir
\textit{heartbeats} a intervalos regulares conmutando entre los estados \texttt{PULSOS\_UP} y
\texttt{PULSOS\_DOWN}. En el control de protección (Nivel 0), otro bloque observa esta señal: si
se detecta una falta de pulsos superando un tiempo preestablecido, el sistema de seguridad asume
una pérdida de comunicación, interrumpe cualquier movimiento en curso y lleva el sistema al
estado de error.

\begin{figure}[H]
    \centering
    \includegraphics[width=0.65\textwidth]{nivel_1/watchdog_heartbeat.png}
    \caption{Watchdog Heartbeat en Stateflow}
    \label{fig:n1_watchdog}
\end{figure}


\subsection{Control del joystick}

Para operar la grúa se utiliza un joystick físico conectado a la estación de operación. Este
dispositivo permite al operador controlar directamente las velocidades de traslación e izaje, así
como activar comandos específicos del sistema. El bloque de manejo de entradas procesa estas
señales para generar las referencias adecuadas para el control supervisor: el eje horizontal del
stick izquierdo se mapea a la velocidad del carro y el eje vertical del stick derecho a la
velocidad de izaje, mientras que los botones gobiernan el cierre de twistlocks y la selección de
modo (encendido/apagado, automático, reset y parada de emergencia). La señal $BTN_{ON/OFF}$
tiene una lógica para funcionar solo con el flanco ascendente del botón, para poder pasar del modo apagado al modo encendido
y no volver a apagado hasta que se pulse nuevamente el botón.  

\begin{figure}[H]
    \centering
    \includegraphics[width=0.85\textwidth]{nivel_1/joystick.png}
    \caption{Bloque de configuración del joystick}
    \label{fig:n1_joystick}
\end{figure}
